%\ifodd\value{page}\hbox{}\newpage\fi
\chapter*{Śrī Lalitā Sahasranāma — Śloka 4}
\addcontentsline{toc}{chapter}{Śloka 4 - Nomi 8,9}

\begin{sloka}
{\sanskritfont
\begin{tabular}{c}
चम्पकाशोकपुन्नागसौगन्धिकलसत्कचा ।\\
कुरुविन्दमणिश्रेणीकनत्कोटीरमण्डिता ॥ ४ ॥
\end{tabular}

\vspace{1.5em}

{\middlesize \iastfont
campakāśoka-punnāga-saugandhika-lasat-kacā |\\
kuruvinda-maṇi-śreṇī-kanat-koṭīra-maṇḍitā ||4||}

\vspace{1em}

{\middlesize
\anvaya{%
चम्पकाशोकपुन्नागसौगन्धिकलसत्कचा
कुरुविन्दमणिश्रेणीकनत्कोटीरमण्डिता ।
}
}

\middlesize \iastfont \itmean{%
«I cui capelli sono profumati dai fiori di \textit{campaka}, \textit{aśoka} e \textit{punnāga};
adornata da un diadema splendente di file di gemme \textit{kuruvinda}.»%
}

}
\end{sloka}

\wordmeaning{%
\wordline{चम्पकाशोकपुन्नागसौगन्धिकलसत्कचा}{campakāśoka-punnāga-saugandhika-\\
lasat-kacā}%
  { dai capelli splendenti e profumati di \textit{campaka}, \textit{aśoka} e \textit{punnāga};
\begin{itemize}
  \setlength{\itemsep}{0pt}
  \setlength{\parskip}{0pt}
  \setlength{\parsep}{0pt}
\item \textit{campaka} (campaka / champak): albero dai fiori fortemente profumati, identificato con \textit{Magnolia champaca} (sin. \textit{Michelia champaca}); fiori giallo‑arancio o bianchi, usati in ghirlande e offerte.
\item \textit{aśoka} (ashoka): “albero aśoka”, \textit{Saraca asoca} (o \textit{Saraca indica}), albero sacro legato a Kāma, al \textit{Rāmāyaṇa} e alla nascita del Buddha, con fiori giallo‑arancio o rosso‑arancio.
\item \textit{punnāga} (punnaga): albero da fiore e da olio, di solito identificato con \textit{Calophyllum inophyllum} (Alexandrian/Indian laurel), sempreverde costiero con fiori bianchi profumati.
\end{itemize}


​  
  }%
\wordline{कुरुविन्द-मणि-श्रेणी-कनत्-कोटीर-मण्डिता}{kuruvinda-maṇi-śreṇī-kanat-\\
koṭīra-maṇḍitā}
  { adornata da un diadema scintillante di file di gemme \textit{kuruvinda}, un tipo di rubino/cristallo rosso, una gemma preziosa dalla forte luce cremisi;}%
}

Lo \textit{Śloka} 4 inaugura la sequenza descrittiva della forma di \textit{Lalitā}, spostando l’attenzione dall’azione cosmica alla manifestazione sensibile della sua bellezza. I profumi floreali non indicano ornamento esteriore, ma la diffusione spontanea della presenza divina, mentre il diadema di gemme esprime la regalità luminosa che incorona la coscienza già descritta nei precedenti \textit{śloka}. La bellezza qui non è decorazione, ma irradiazione visibile dell’ordine e della pienezza del principio divino.

Lo \textit{Śloka} 4 presenta due nāma effettivi, entrambi appartenenti alla sequenza descrittiva della forma sensibile di \textit{Lalitā}. \textbf{\textit{Campakāśoka-punnāga-saugandhika-lasat-kacā}} definisce la Dea come colei i cui capelli sono permeati del profumo dei fiori di \textit{campaka}, \textit{aśoka} e \textit{punnāga}, indicando una bellezza che non si impone visivamente, ma si diffonde come irradiazione naturale della presenza divina. \textbf{\textit{Kuruvinda-maṇi-śreṇī-kanat-koṭīra-maṇḍitā}} afferma invece la sua regalità luminosa, espressa dal diadema ornato di file di gemme \textit{kuruvinda} scintillanti. In questo \textit{śloka} i \textit{nāma} non svolgono una funzione operativa o cosmica, ma introducono la contemplazione diretta della forma, mostrando come la pienezza della coscienza si renda visibile e percepibile attraverso la bellezza ordinata del corpo divino.
